% Generated draft: requires mathematical and source review.
\documentclass[11pt]{article}
\usepackage[margin=25mm]{geometry}
\usepackage{fontspec}
\setmainfont{Latin Modern Roman}
\usepackage{amsmath,amssymb,mathtools,unicode-math}
\setmathfont{Latin Modern Math}
\usepackage{enumitem,needspace,xurl,hyperref}
\setlength{\parindent}{0pt}
\setlength{\parskip}{5pt}
\setlength{\emergencystretch}{3em}
\newcommand{\N}{\mathbb{N}}
\newcommand{\Z}{\mathbb{Z}}
\newcommand{\Q}{\mathbb{Q}}
\newcommand{\R}{\mathbb{R}}
\newcommand{\C}{\mathbb{C}}
\newcommand{\F}{\mathbb{F}}
\newcommand{\A}{\mathbb{A}}
\newcommand{\PP}{\mathbb{P}}
\newcommand{\E}{\mathbb{E}}
\newcommand{\eps}{\varepsilon}
\newcommand{\dd}{\,\mathrm{d}}
\begin{document}
\section*{UnsolvedMath ID: 30004303}
\textbf{Source problem code:} catalogue entry
\subsection{Definitions and mathematical statement}
For a set \(X\), a map \(r\colon X\times X\to X\times X\) satisfies the set-theoretic Yang--Baxter equation in braid form if
\[
(r\times\operatorname{id}_X)(\operatorname{id}_X\times r)(r\times\operatorname{id}_X)
=
(\operatorname{id}_X\times r)(r\times\operatorname{id}_X)(\operatorname{id}_X\times r)
\quad\text{on }X^3.
\]
Under the usual structure-monoid convention, define
\[
M(X,r)=\left\langle X\;\middle|\;xy=uv\text{ whenever }x,y,u,v\in X\text{ and }r(x,y)=(u,v)\right\rangle.
\]
The structure algebra \(A=kM(X,r)\) is the monoid algebra: its elements are finite formal sums \(\sum_{m\in M(X,r)}a_m m\), with coefficients in \(k\), and multiplication extends the monoid multiplication distributively, so that \((a m)(b n)=(ab)(mn)\). This description presupposes an associative coefficient ring with identity; whether the intended coefficient domain is specifically a field, and what characteristic restrictions apply, are not supplied.

A ring \(A\) is \emph{semiprime} if every two-sided ideal \(I\) of \(A\) satisfies
\[
I^2=\{0\}\ \Longrightarrow\ I=\{0\}.
\]

\textbf{Problem.} Determine necessary and sufficient conditions on the coefficient domain \(k\) and the Yang--Baxter solution \((X,r)\) for \(kM(X,r)\) to be semiprime, and determine what structural restrictions on \((X,r)\) follow from this semiprimeness.

\textbf{Unresolved scope.} This is a characterization request, not a specified conjectural equivalence. The supplied record does not fix the admissible coefficient domains, whether \(X\) is finite, or whether solutions must be bijective, nondegenerate, involutive, or idempotent. Some conventions include bijectivity in the definition of a set-theoretic solution; the equation above alone does not impose it. The definitions above explain the standard notation without resolving these missing hypotheses. The intended range of quantification and the particular structural consequences sought require clarification from the source before a fully scoped statement can be certified.
\subsection{Short English statement}
Characterize when the structure algebra \(kM(X,r)\) is semiprime, meaning that it has no nonzero two-sided ideal whose square is zero, and determine what this implies about the Yang--Baxter solution \((X,r)\). The supplied record leaves the coefficient domain and the class of solutions unspecified, so this remains a statement-only draft requiring mathematical review.
\subsection{Sources}
\begin{itemize}
\item Supplied local UnsolvedMath registry record: permanent ID \(30004303\), title \emph{Semiprime Yang--Baxter Structure Algebras}, exact source problem code \texttt{catalogue entry}, in the Oberwolfach Reports Open Problems collection. The record gives proposed year \(2019\) and names no individual proposer.
\item Supplied report citation: \emph{Mini-Workshop: Algebraic Tools for Solving the Yang--Baxter Equation} (2020), Oberwolfach Reports, pp.~3207--3242, DOI \texttt{10.4171/owr/2019/51}. \url{https://doi.org/10.4171/owr/2019/51}. The background attributes the problem collection to workshop open-problem discussions. No problem-specific page or speaker is supplied.
\item Supplied literature reference: Colazzo--Jespers--Kubat--Van Antwerpen, structure algebras. \url{https://arxiv.org/abs/2305.06023}.
\item Supplied literature reference: Colazzo et al., finite idempotent solutions. \url{https://academic.oup.com/imrn/article/2024/7/5458/7240113}.
\item Supplied literature reference: 2026 structure-algebra results. \url{https://doi.org/10.1016/j.jalgebra.2026.04.044}.
\item Literature assessment dated \texttt{2026-08-22}: the supplied curation labels the problem partially solved and reports semiprimeness characterizations in several finite bijective and idempotent families, together with recent structural consequences, while reporting no classification of its effect on arbitrary Yang--Baxter solutions. This is supplied dated curation derived from the supplied OWR material, not an independently verified literature assessment; the status and cited claims require independent verification.
\end{itemize}
\end{document}
